\scanpage{149}
\deleted{\unsure{indique que|implique que}} Le composé de (\unsure{22|21}) avec \deleted{$P^\infty_{S/T}$}
\[
P^\infty_{S/T}(E)^+\longrightarrow P^1_{S/T}(E)^+\simeq E\otimes\Omega^1_{S/T},\ \deleted{\text{\unsure{est}}}
\]
noté $d_E^0$
\[
d_E^0:E\longrightarrow E\otimes\Omega^1_{S/T},\quad d_E^0=d_E^+\bmod\ldots \tag*{(23)}
\]
est l'hom (5) associé à la connexion de $E$ définie
par la donnée de recollement. \deleted{Si} \uline{On vérifie que si
la donnée de recollement} \deleted{définie par $\delta_E$} \uline{est une stratification},
alors la connexion associée est intégrable, i.e. (23)
donne lieu à un complexe de De Rham.

On sait \unsure{ds|alors|bien} [\quad] définir des hom canoniques
\[
H^*_{\mathrm{strat}}(S,E)\longrightarrow H^*_{\mathrm{DR}}(S,E), \tag*{(24)}
\]
$\unsure{\partial|\delta}$-fonctoriels en $F$, compatibles au changemt de base.

Nous voulons une interprétation de $H^1_{\mathrm{strat}}(S,E)$
analogue à celle donnée plus haut pour $H^1_{\mathrm{DR}}(S,E)$.

\deleted{Pour ceci, \deleted{\unsure{posons}} Soit, \inserted{pour tt $\unsure{i\in\mathbb N|i\geq0}$,} $C^i(E)$
\deleted{l'image inverse de $E$} \deleted{\unsure{sur le}} le faisceau sur
$S_{\mathrm{Zar}}$ \deleted{image inverse du} complété formel le long de la diagonale de
$(S/T)^{i+1}$ de l'image inverse de $E$ par (disons) $\operatorname{pr}_{\unsure{1|i}}$.
\deleted{La stratification de}}

\[
\deleted{C^0(E)=E,\quad C^1(E)=P^\infty_{S/T}(E).}\tag*{\deleted{(24)}}
\]

\deleted{La stratification de $E$ met sur $C^{\unsure{\cdot|\bullet}}(E)$ une structure
de faisceau cosimplicial [en fait, un co-module
cosimplicial sur $C^{\cdot}(\underline O_S)$] :}

\diagram{\begin{tikzcd}[ampersand replacement=\&,column sep=small,row sep=small,execute at end picture={\draw[thin] (current bounding box.south west) -- (current bounding box.north east);\draw[thin] (current bounding box.north west) -- (current bounding box.south east);}]
\deleted{(25)}\&\deleted{\deleted{\unsure{E^0|E}}}\arrow[r,shift left,"\deleted{\delta_E^\infty}"]\arrow[r,shift right,swap,"\deleted{d_E^\infty}"]\&
\deleted{\deleted{P_{S/T}^\infty(E)}}\&\deleted{C^0(E)}\arrow[r,shift left]\arrow[r,shift right]\&
\deleted{C^1(E)}\arrow[r,shift left=1.2ex]\arrow[r]\arrow[r,shift right=1.2ex]\&\deleted{C^2(E)\ldots}\\
\&\&\&\deleted{E}\arrow[u,equal]\arrow[r,shift left,"\deleted{\delta_E^\infty}"]\arrow[r,shift right,swap,"\deleted{d_E^\infty}"]\&
\deleted{P^\infty_{S/T}(E)}\&\deleted{,}
\end{tikzcd}}
% archive sheet number: 148
